\documentclass[11pt]{article}
\usepackage[utf8]{inputenc}
\usepackage{graphicx,amsmath}
\usepackage{booktabs}
\usepackage{multirow}
\usepackage[margin=1in]{geometry}
\usepackage{hyperref}
\title{What deleting an agent's memory actually removes: evidence from three preregistered iterations}
\author{Liu Zewen\thanks{ORCID 0009-0003-2981-9888; GitHub: \texttt{github.com/aidless}}}
\date{20 November 2026}
\begin{document}
\begin{sloppypar}
\maketitle

\begin{abstract}
Deleting a stale memory seems useful, yet three preregistered iterations show each apparent benefit of deletion traces to a different source. A two-arm contrast appeared to confirm that removing a plan-attribution entry cuts stale-plan following; a factorial decomposition showed the prompt carried the whole effect while deleting the entry was inert. A conversation-history carrier made the arms differ in one visible object: deletion produced 119 of 119 stale-write episodes at 3.2--3.8 times the token cost. Resource identifiers existed only inside the deleted carrier, so deletion landed no accepted write: the manipulation removed the vocabulary for acting, not a misleading belief. The third iteration fixes those identifiers in the system prompt, leaving only a version mapping in the carrier. With that confound removed, deletion still raises stale writes (2 of 159 versus 53 of 159, +32.1 percentage points, McNemar $p=1.2\times10^{-14}$) while delivery is statistically indistinguishable (199 of 200 versus 200 of 200); deletion writes 1481 times before invalidation, and the token ratio falls to 1.0177. A carrier-sensitivity replication moves that mapping into a retrieval corpus with identifiers fixed; there deletion again raises stale writes (79 of 146 effective pairs versus 10, +47.3 points), delivery falls for the first time (135 of 147 versus 147 of 147), at a token ratio of 1.2622. In the three iterations deletion cost the task in none: the apparent effects were prompt wording, missing identifiers, or a corrected stale write. The replication is the exception: with the version tag reachable only via retrieval, deletion cost delivery in 12 of 147 episodes. Deletion removes whatever else the carrier held; the cost depends on re-derivability, not on staleness.
\end{abstract}

\section{Introduction}

Agent memory is usually evaluated as an accumulator: more retained context is assumed to be better, and forgetting is studied as a budget question \cite{wang2023augmenting,xu2025amem,hu2025evaluatingmemory,packer2023memgpt,zhong2023memorybank,park2023generative}. Architectural accounts distinguish what is stored from where it is stored and from what the agent can re-derive \cite{sumers2023cognitive,zhou2026externalization}. A different failure mode appears when memory is \emph{wrong} rather than missing. The agent holds a plan, the environment changes underneath it, and the stored plan is now the reason the agent keeps failing \cite{chao2026stale}; this is an instance of context--memory conflict in which stored and current information disagree \cite{xu2024knowledge}. The natural remedy --- delete the offending entry --- has an intuitive appeal: if the plan is the problem, remove the plan \cite{wei2026fademem}. A recent benchmark asks whether an agent complies when asked to forget a stored fact \cite{xie2026unlearning}; the question here is different --- what deleting one stored entry changes in behaviour when the fact survives elsewhere or must be re-derived, and which channels mediate that change. The proposition under test is that deleting an entry changes behaviour only when that entry is the unique available carrier of decision-relevant information; where another reachable channel holds the same fact, deletion should cost at most the work of re-deriving it.

This paper asks whether that remedy works, under a condition we designed to be as favourable to it as possible. We constructed a task where a stored plan does become stale, where the environment rejects writes that follow it, and where the agent is told that a change occurred. We then varied exactly one thing: whether a durable record that can carry the plan forward exists in the agent's own history. If deleting the culprit entry helps because it removes a misleading belief, then the deletion arm should recover. If it does not help --- because the deleted record carries more than a plan --- then the deletion arm should keep failing, and it should fail expensively. We preregistered the outcome variable as stale-plan following; the experiment returns a sharper result than that variable can express, and we report both the preregistered statistic and the diagnostic that disambiguates it.

The study is the second iteration of a preregistered line. In the first iteration we named a specific mechanism, ``removing the attribution entry that records the plan recovery path reduces stale-plan following'', and ran a packaged two-arm comparison that appeared to confirm it: the arm with the entry removed logged a stale write in 0 of 40 paired episodes against 40 of 40 in the control, at a token cost ratio of 1.1479. A $3\times2$ factorial decomposition then showed that the packaged contrast was confounded, consistent with the documented sensitivity of model behaviour to prompt wording and example order \cite{sclar2023sensitivity,zhao2021calibrate}. Holding memory fixed and changing only the instruction prompt moved the stale rate from 5/6 to 0/6 cells; holding the prompt fixed and swapping the culprit entry for an unrelated entry moved almost nothing (193 paired episodes in the authority cells, difference 1.0 percentage point; 83 and 107 paired episodes in the recheck cells, difference 0.0). The packed effect was real; the named mechanism was not what produced it. We therefore did not scale the first design. Adding episodes to a confounded contrast buys precision on the wrong quantity.

The present study is the redesign that the decomposition implied. The instruction prompt is frozen to a single constant in both arms. The carrier is not an instruction at all: it is two ordinary messages in the conversation history that record the migration. Deleting them changes what the agent has seen, not what it has been told to prioritise. Everything else that could leak the answer is removed or made identical, and we verify that removal by reading the returned tool payloads rather than by assuming it.

Our contributions are: (i) a carrier-based operationalisation of memory deletion in which the two arms differ in exactly one visible object, with the identifying channels audited; (ii) three successive preregistered iterations whose apparent deletion effects each dissolve into a different mechanism when the right diagnostic is read --- instruction wording in the first, identifier availability in the second, and the counting of a corrected stale write in the third; (iii) a carrier-sensitivity replication that moves the surviving version mapping from conversation history into a retrieval corpus and finds the stale-write effect persist (+47.3 points over effective pairs) while delivery, for the first time, degrades; and (iv) an honest account of those mechanism failures and of the measurement artifacts (a 400-token output cap that turned truncation into apparent task failure) that the redesign had to fix before any of the present numbers could be trusted. The through-line is a negative result about interpretation rather than about deletion: across all three iterations, deleting the entry did not make the agent fail; each failure was something else wearing deletion's clothes.

\section{Methods}

\subsection{Task and environment}

Each episode is a 40-file delivery task, a setting chosen to expose plan following under tool use \cite{qin2023toolllm,li2025sopbench,yehudai2025survey}. The agent must write 40 files to a target prefix under a versioned path scheme. The environment exposes two stores, \texttt{store-a} and \texttt{store-b}, and a target of the form \texttt{/dst/<seed>/v1}. After the agent has written a small number of files, the environment mutates: the writable target becomes \texttt{/dst/<seed>/v2}, and writes to the \texttt{v1} prefix are rejected with an error whose text says the stored plan no longer matches the environment. The task is complete when all 40 files are verified at the live target. Cards are generated deterministically with 40 files each and a per-card seed; the same 40 cards are used in both arms and in all three model batches.

\subsection{Carrier manipulation}

Externalization distinguishes the store in which a fact is written from the content that makes it actionable \cite{zhou2026externalization,sumers2023cognitive}. In this study the carrier (condition S\_c) is delivered as two messages appended to the conversation history: one user-role message containing a two-line migration memo that maps the old prefix to the new one, and one assistant-role acknowledgement. In the \emph{keep} arm these two messages remain in history for the whole episode. In the \emph{del} arm they are removed before the first turn. Nothing else differs: the system prompt is a single constant string shared by both arms, the tool set is identical, the invalidation notice injected after the change is the same sentence in both arms, temperature is 0 and the output cap is identical.

\subsection{Identification control}

Two design choices prevent the instrument from answering the question for the agent. First, the \texttt{stat} tool returns a store status record that contains the fact of the change but no target, no resource and no plan, so it does not serve as a substitute carrier. Second, when a write is rejected the tool returns an error code and a hint that the plan is outdated, but it does not return the correct target or resource. An agent can only recover by re-deriving the target itself. This matters because an earlier frozen version of the instrument broadcast the invalidation signal in a form that made recovery free; we treat a free signal as a design defect, not as a finding about memory.

\subsection{Verdicts}

Each episode is assigned one of five verdicts by a deterministic judge over the recorded trace: RECOVERED (all files verified at the live target), STALE\_FAILURE (the frozen judge's label for at least one write issued against a superseded version; we report the quantity it counts as stale-write incidence and reserve ``failure'' for the delivery endpoint), INCOMPLETE (neither, with a non-degenerate trajectory), INCOMPLETE\_TRUNCATED (the final turn was cut off by the output cap and produced no tool call), and PRE\_INJECTION\_COMPLETION (all files verified before the invalidation event fired). The last two exist because the first iteration conflated them with task failure. Truncation is a property of the instrument; pre-injection completion is a property of the arm's speed and does not occur in an arm that does not receive the carrier to begin with. Neither enters the paired denominator. We report them separately rather than folding them into either outcome.

\subsection{Preregistration and frozen code}

The second-iteration protocol is in force from 17 November 2026; an earlier draft fixed a cost gate of 1.10 that the final protocol dropped, and two of the three reported batches had completed before the final protocol took effect. It fixes the carrier, the arms, the verdict set, the statistic (paired exact McNemar on discordant pairs, with the paired unit being a card), the cost definition (mean per-episode tokens in the deletion arm divided by the same in the retention arm, over effective pairs), the transport-loss rule (a batch in which more than 20\% of episodes failed at the transport layer is void; a rate above 50\% demotes it to directional evidence), and the requirement that at least one model have zero transport loss. It does not presuppose the direction of the cost effect: the first iteration found deletion cheaper because it failed fast, a reason that does not transfer to a design in which deletion fails slowly, so we report this iteration's cost direction as exploratory.

The six files that constitute the measurement instrument are frozen by hash and are byte-identical to the first iteration's registration; a provenance block recomputes their hashes on demand and a gate exits non-zero on mismatch. Orchestration code is versioned separately and every change to it is recorded in a deviation log with the reason, the reproduction command and the observed effect. All deviations from the first iteration's plan --- twenty-eight of them by the time of writing --- are listed in \texttt{DEVIATION\_L3.md}; the ones that bear on the present numbers are the carrier redesign, the output-cap change, the truncation verdict, and the pre-injection verdict.

\subsection{Models and execution}

Three models were run through an OpenAI-compatible gateway:\newline
\noindent\path{openrouter/minimax/minimax-m3.1-flash}, \path{moonshotai/kimi-k2-thinking}, and \path{workbuddy/deepseek-v4.1-flash}. Each batch runs 40 cards in both arms, i.e.\ 80 episodes, at two to three concurrent workers. Token accounting is the sum of prompt and completion tokens reported by the provider for the episode. Episodes whose failure is at the transport layer (HTTP errors that exhausted retries) are marked, excluded from all cells, and removable by a purge tool that keys on the transport flag rather than on the verdict, so that a genuine failure is not recorded as a transport loss.


The output cap is 1600 tokens, above the reasoning-budget failure range reported for comparable models \cite{pu2025thoughtterminator,huang2025adactrl}. This is a deviation from the first iteration's 400 and it is load-bearing: with a 400-token cap on a reasoning model, 48 of 270 recorded turns (17.8\%) hit the cap, 12 of those produced no tool call, and the frozen harness treats a tool-call-free turn as the model choosing to stop. In a controlled reproduction, re-sending the same context with a 1600-token cap produced tool calls in 12 of 12 such turns. The apparent incompleteness in the first iteration was therefore substantially an artifact of the cap.

\subsection{Positive control}

Before each batch is interpreted, an 8-card positive control is run on a model that is known to expose the contrast: the retention arm must recover in at least 80\% of cards and the deletion arm must degrade by at least 50 percentage points. The control for the reported batches used \texttt{workbuddy/deepseek-v4.1-flash}: retention recovered 8/8 (median 14080 tokens), deletion completed 0/8 (median 105438 tokens). Reading the control traces shows why the deletion arm completes nothing: every \texttt{write} it issues names a resource that does not exist, and the service rejects each with \texttt{resource\_unavailable}. The control therefore confirms that the manipulation is large, but it also shows that the manipulation removes the identifiers, not merely the plan. A control that fails to show this contrast voids the batch rather than being reported as a negative result.

\subsection{Statistics}

The primary statistic is the paired difference in stale-write incidence between the deletion and retention arms, over cards in which both arms produced an effective verdict, with an exact two-sided McNemar test on the discordant counts \cite{peyrard2021paired,yang2023mcnemar}. The cost statistic is the ratio of mean per-episode tokens, reported with a paired bootstrap interval (10000 resamples, seed 20261117). The preregistered sample target for the cost interval was 420 effective pairs; the three batches together yield 109, and we report the shortfall rather than treating the achieved sample as if it were the target.

\section{Results}

\subsection{Main contrast}

Table 1 gives the per-batch cell counts. The deletion arm produced no recovered episode in any batch. The retention arm produced one stale-write episode across 110 effective episodes: card \texttt{L3V2SC-216000}, which wrote against the superseded version once and then went on to verify all 40 files at the live target --- the compound outcome that the first iteration's single-label judge reported as a pure failure, and which we now report as what it is.

\begin{table}[t]
\centering
\caption{Per-batch cell counts over 40 cards per arm. Effective units are RECOVERED plus STALE\_FAILURE; transport losses are excluded from all cells; truncated and pre-injection completions are shown separately and do not enter the paired denominator.}
\label{tab:1}
\begin{tabular}{llrrrrrrr}
\toprule
Batch & Arm & $n$ & TL & REC & STALE & INC & TRUNC & PIC \\
\midrule
\multirow{2}{*}{minimax} & keep & 40 & 0 & 36 & 0 & 1 & 0 & 3 \\
 & del  & 40 & 0 & 0 & 40 & 0 & 0 & 0 \\
\midrule
\multirow{2}{*}{kimi} & keep & 38 & 2 & 34 & 0 & 0 & 0 & 4 \\
 & del  & 39 & 1 & 0 & 39 & 0 & 0 & 0 \\
\midrule
\multirow{2}{*}{workbuddy} & keep & 40 & 0 & 39 & 1 & 0 & 0 & 0 \\
 & del  & 40 & 0 & 0 & 40 & 0 & 0 & 0 \\
\bottomrule
\end{tabular}

\end{table}

Table 2 gives the paired statistics. The per-model difference is between 97.5 and 100 percentage points, with all discordant pairs moving in the same direction, and the exact McNemar $p$ is $2.9\times10^{-11}$, $2.3\times10^{-10}$ and $3.6\times10^{-12}$ in the three batches, respectively.

\begin{table}[t]
\centering
\caption{Paired contrasts, with the deletion-minus-retention token ratio and its paired bootstrap interval. The last column is the preregistered cost gate from the first iteration, shown for continuity only; the second-iteration protocol does not presuppose a cost direction.}
\label{tab:2}
\begin{tabular}{lrrrrrr}
\toprule
Batch & pairs & keep stale & del stale & $\Delta$ (pp) & McNemar $p$ & cost ratio [95\% CI] \\
\midrule
minimax & 36 & 0.000 & 1.000 & 100.0 & $2.91\times10^{-11}$ & 3.21 [2.78, 3.78] \\
kimi & 33 & 0.000 & 1.000 & 100.0 & $2.33\times10^{-10}$ & 3.82 [3.00, 4.86] \\
workbuddy & 40 & 0.025 & 1.000 & 97.5 & $3.64\times10^{-12}$ & 3.66 [3.15, 4.23] \\
\midrule
pooled & 109 & --- & --- & --- & --- & 3.49 [3.17, 3.86] \\
\bottomrule
\end{tabular}

\end{table}

\subsection{What the deletion arm actually did}

The four verdict categories record outcomes, not activity, and the activity is what disambiguates the result. Table 3 reports how many \texttt{write} calls the environment accepted. The deletion arm landed no successful write in any of the three batches --- not after the invalidation event and not before it --- while the retention arm landed writes from the first or second step in 109 of 110 effective episodes. The resource identifiers the task requires appear only inside the deleted carrier, so once that carrier is removed the deletion arm spends its budget proposing resource names that the service does not know and rejects with \texttt{resource\_unavailable}. The failure mode is thus an absent identifier rather than a followed stale plan. The paired contrast measures the removal of information that the environment will not restate, which is a stronger manipulation than the preregistered variable of stale-plan following names, and we report both the preregistered statistic and the diagnostic that separates them.

\begin{table}[t]
\centering
\caption{Write-level diagnostic per batch and arm. \emph{pre} and \emph{post} are accepted over attempted \texttt{write} calls before and after the invalidation event; \emph{zero-write} counts episodes with no accepted write, out of the arm size.}
\label{tab:3}
\begin{tabular}{llrrrr}
\toprule
Batch & Arm & pre ok/total & post ok/total & never ok & $n$ \\
\midrule
\multirow{2}{*}{minimax} & keep & 346/357 & 248/383 & 1/40 & 40 \\
 & del  & 0/224 & 0/3322 & 40/40 & 40 \\
\midrule
\multirow{2}{*}{kimi} & keep & 350/465 & 228/228 & 1/40 & 40 \\
 & del  & 0/334 & 0/1765 & 40/40 & 40 \\
\midrule
\multirow{2}{*}{workbuddy} & keep & 156/168 & 193/253 & 0/40 & 40 \\
 & del  & 0/104 & 0/2385 & 40/40 & 40 \\
\bottomrule
\end{tabular}

\end{table}

\subsection{Cost}

Deleting the carrier is more expensive, not cheaper: the pooled ratio is 3.49 with a 95\% bootstrap interval of [3.17, 3.86], and every batch interval lies above both 1.0 and the 1.10 threshold that the first iteration used as its gate. The mechanism is visible in the medians. Retention-arm episodes finish in a median of 16{,}476 tokens; deletion-arm episodes spend a median of 64{,}869 tokens before being judged, because the agent keeps re-issuing writes that the environment keeps rejecting, and re-planning from a state that no longer contains the mapping. In the first iteration deletion looked cheap for the opposite reason: there the deletion arm abandoned early. Failure speed, not failure rate, is what the first cost ratio measured.

\subsection{Pre-injection completion asymmetry}

Four verdict categories do not express an asymmetry that the design makes structural. An episode can complete the whole task before the invalidation event fires, and such an episode has no stale write because it did not meet a stale state. This can only happen in the arm that is fast enough to finish early --- in practice, the retention arm. We observed 7 such episodes (3 minimax, 4 kimi, 0 workbuddy), all in the retention arm, and 0 in the deletion arm across all three batches. Reporting them as incompleteness would understate retention performance; reporting them as recovery would overstate it. They are therefore listed separately and excluded from the paired denominator, and we take the asymmetry itself as a design property rather than as evidence for the hypothesis.

\subsection{What the first iteration's mechanism failure implies}

The first iteration reported a large packaged effect and named a mechanism for it. The factorial decomposition that followed is what makes the present design interpretable, so we summarise its outcome. Holding the prompt at the authority variant and swapping the deleted entry between the culprit and an unrelated entry changed the paired difference by 1.0 percentage point over 193 episodes, and by nothing at all over 83 and 107 episodes in the recheck cells. Holding memory fixed and changing the prompt from authority to recheck moved the stale-write incidence from 5 of 6 cells to 0 of 6, and in the larger batch produced a paired difference of 99.0 percentage points over 105 episodes, which on 105 pairs can only be 104 discordant episodes in one direction and none in the other (exact two-sided $p<10^{-30}$). Memory deletion was inert; the instruction prompt carried the entire packed effect. The present study keeps the prompt constant precisely because of that result, and the result it obtains --- a stable 97.5--100 point gap with the prompt held fixed --- is therefore about the carrier, not about the words the agent was given.

\subsection{A third iteration: keeping the identifiers, removing only the version}
\label{sec:v3}

The second iteration's result has a confound that its own activity diagnostic exposes. The resource identifiers (\texttt{store-a}, \texttt{store-b}) appear only inside the deleted carrier; the system prompt does not contain them. Deleting the carrier therefore removes both the plan's version and the vocabulary needed to act on it, and the deletion arm lands no accepted write at all. That is a stronger manipulation than stale-plan following names, and it does not tell us what happens when only the stale version is removed.

We ran a third, separately preregistered iteration to isolate that case. The resource identifiers, the target shape and the two stores are now stated permanently in the system prompt, byte-identical in both arms. What the carrier carries is only the version mapping --- that the live prefix has moved from version 1 to version 2. The \emph{keep} arm retains this mapping; the \emph{del} arm loses it. Both arms can write from the first turn: in the deletion arm, 1481 writes are accepted before the invalidation event fires, and both arms land at least one accepted write in all 200 episodes. The manipulation now removes only the knowledge of which version is current. The registered hypothesis predicted that removing the version would reduce stale writes by at least 15 points with the direction holding in each model; the observed effect went the other way (+32.1 points) and was measured on one model, so we report it as exploratory rather than as confirmation.

Under this clean manipulation the paired stale-write effect survives but shrinks, and its meaning changes. Table 4 gives the cell counts. The deletion arm produces stale writes in 53 of 159 effective pairs versus 2 for retention, a difference of +32.1 percentage points (exact McNemar $p=1.2\times10^{-14}$, well above the registered 140-pair minimum). But the stale write is not a failure to deliver: all 59 deletion-arm stale episodes verify their 40 files at the live target, including the corrected switch after the first rejection. Delivery rates are 199/200 in retention and 200/200 in deletion --- statistically indistinguishable --- and the token ratio collapses from 3.49 in the second iteration to 1.0177, inside the 1.10 cost gate. The 32-point gap therefore measures whether the agent ever issued a write against the superseded version, not whether it completed the task. Deletion's cost in this regime is a brief detour through a rejected write, not a failure to finish.

\begin{table}[t]
\centering
\caption{Third iteration, 200 cards per arm. \emph{delivered} counts episodes verifying all 40 files at the live target; \emph{ok before inj.} counts accepted \texttt{write} calls before invalidation. The final row is the effective paired contrast (deletion stale vs.\ retention stale) and the paired token ratio.}
\label{tab:4}
\begin{tabular}{lrrrrrr}
\toprule
Arm & $n$ & REC & STALE & delivered & ok before inj. & median tok. \\
\midrule
keep & 200 & 174 & 2 & 199/200 & 1276 & 15,737 \\
del  & 200 & 123 & 59 & 200/200 & 1481 & 16,086 \\
\midrule
paired & 159 & --- & 53 vs 2 & --- & --- & $\times$1.0177 \\
\bottomrule
\end{tabular}

\end{table}

\subsection{A carrier-sensitivity replication}
\label{sec:sa}

The third iteration tests one carrier, conversation history. If the remaining 32-point effect is about the version mapping rather than about that particular carrier, moving the mapping to a different carrier should keep the effect and should not disturb delivery. We therefore ran a carrier-sensitivity replication: the version mapping lives in a retrieval corpus that the agent can query with a \texttt{retrieve} tool, an instance of retrieval-augmented generation in which a non-parametric store supplies the fact \cite{lewis2020rag}, the resource identifiers and target shape remain in the system prompt exactly as in the third iteration, and the \emph{del} arm has the version lines stripped from the corpus while keeping the task description. This arm was designed after the third iteration and is reported as a sensitivity result, not as a fourth preregistered iteration. The instrument is a separate runner with its own frozen-judge reuse; run on a single model (\texttt{workbuddy/deepseek-v4.1-flash}) over 147 cards per arm, with 12 of 306 episodes lost at the transport layer and excluded.

Under this carrier the effect survives and delivery does not. Table 5 gives the cell counts. The deletion arm produces stale writes in 79 of 146 effective pairs versus 10 for retention, a difference of +47.3 percentage points (exact McNemar $p=3.7\times10^{-18}$), larger than the 32 points of the third iteration. Delivery is 147 of 147 in retention against 135 of 147 in deletion --- the first iteration in which deletion costs the task itself --- and the token ratio is 1.2622, above the 1.10 gate.

The larger gap has a second source that the cell counts hide. Fifty of the 147 deletion-arm episodes land no accepted write before the invalidation event at all; they spend the episode repeatedly querying the corpus and proposing version-0 writes that the environment rejects, which is a different behaviour from following a stale version. Because the manipulation itself determines whether an episode enters that stratum, conditioning on it is a selection rather than a controlled comparison, so no stratified estimate is reported in the main text; the subgroup counts appear in Appendix~\ref{app:sa-strat} for description only. What the replication adds is not a larger effect but a cost that the earlier carriers did not have: with the mapping in a retrieval corpus, deletion can consume the episode rather than only costing a corrected mis-step.

\begin{table}[t]
\centering
\caption{Carrier-sensitivity replication (retrieval corpus, 147 cards per arm). \emph{delivered} counts episodes verifying all 40 files at the live target; \emph{ok before inj.} counts accepted \texttt{write} calls before invalidation. The \emph{paired} row uses all effective pairs; the post-treatment subgroup is reported descriptively in Appendix~\ref{app:sa-strat}.}
\label{tab:5}
\begin{tabular}{lrrrrrr}
\toprule
Arm & $n$ & REC & STALE & delivered & ok before inj. & median tok. \\
\midrule
keep & 147 & 137 & 10 & 147/147 & 285 & 36,401 \\
del  & 147 & 67 & 79 & 135/147 & 148 & 50,079 \\
\midrule
paired & 146 & --- & 79 vs 10 & --- & --- & $\times$1.2622 \\
\bottomrule
\end{tabular}

\end{table}

\section{Discussion}

Across the three iterations, the recurring surprise is that the deletion effect keeps changing size and keeps refusing to be about deletion. The first iteration's 100-point gap tracked the wording of the instruction prompt. The second iteration's gap tracked whether the resource identifiers survived the deletion. The third iteration's 32-point gap tracks whether a stale version is present to be followed --- and even there the episodes carrying the stale-write label deliver the task. Each iteration removed one more rival explanation, and in each case what remained was smaller and less about forgetting than the headline number suggested. We report the three together because the trajectory is the result: a single labile number, three different sources.

This reframes what ``memory deletion'' is for. If a stale entry is the problem, deleting it should remove a bad belief and the agent should improve. If the entry is the last remaining record of a fact the environment will not restate, deleting it removes the route back while leaving competence intact --- the second iteration. If the fact is stated elsewhere and only its current version lives in memory, deleting it costs a corrected mis-step but not the task --- the third. These situations look identical from inside the memory store and differ only in what else carries the fact and what the task actually requires. Our identification controls try to make that distinction operational by removing or fixing each alternative channel in turn and measuring what deletion costs in each regime.

Read as a consistency problem, the three regimes are the cases a cache protocol names --- a single authoritative copy, a replica with a reachable authority, and a replica whose authority is expensive or absent --- and they are an instance of the disagreement between stored and current information that motivates work on context--memory conflict \cite{xu2024knowledge}. The analogy stops where the memory store differs from a cache: the copies are natural-language text rather than addressed bytes, their equivalence is judged semantically, and no counter tracks how many replicas exist, so a cache-style invalidate-all-replicas step has no clean implementation here. What a cache protocol settles structurally, an agent has to settle by probing which channels still hold the fact, and that is why the cost of deletion, not its direction, is what varies across our regimes. The operational form of the proposition is a pre-deletion audit: enumerate the channels that can carry the fact, probe whether any restates it and at what cost, and apply a three-way rule --- retain an entry that is the only carrier of something the task needs, delete one that another channel restates cheaply, and externalise a pointer when the other channel exists but is expensive. The audit is instrumentation a deployment must perform for itself.

The carrier-sensitivity replication tests whether that account survives a change of carrier, and it partly does. Moving the version mapping from conversation history into a retrieval corpus keeps the stale-write effect at the same order of magnitude, so the effect is not an artifact of the history carrier. What changes is the cost: in the retrieval corpus the deletion arm can spend its budget querying and proposing writes that the environment rejects, and 12 of 147 episodes end without delivering the task. This is the one regime we observed in which deletion costs the task rather than a corrected mis-step, and it occurs when re-derivation is possible but costly.

The cost numbers make the same point from a different angle. A deletion policy judged on stale-write incidence alone would look badly harmful in the second iteration (a near-100-point swing) while hiding a 3.5-fold cost multiplier; in the third iteration the same rate gap of 32 points comes with a token ratio near 1.0; in the carrier-sensitivity replication the gap widens to 47 points and delivery drops for the first time, at a ratio of 1.2622. The first iteration's ratio of 1.1479 made deletion look nearly free for a third reason: there the deletion arm abandoned early. These are correct measurements of different regimes, and none is a property of ``deletion'' as such. This is why we report each token ratio with its interval and refuse to carry a direction prior into the design.

\section{Limitations}

The second iteration's preregistered cost target was 420 effective pairs; its three batches achieved 109, and the third iteration cleared its own 140-pair minimum with 159, so the reported intervals are honest about which registered design each belongs to. The point estimates are stable and the intervals exclude the thresholds of interest, but the second-iteration intervals are wider than its target design would have produced. The shortfall is a channel-availability artifact, not a design choice, and it may understate how wide the intervals would be at the registered sample: the gateway's model endpoints were rate-limited in windows that lasted hours, and one model that had failed completely came back the same day. Batches were therefore assembled one at a time, with a transport-loss purge between rounds, which is also why one batch's report file records a token ratio (3.7622) that differs by 1.4\% from the value recomputed from its own judgements (3.8151); the paper cites the recomputed value and the discrepancy is documented. Third, two carriers are tested: S\_c (conversation history) and, in the sensitivity replication, a retrieval corpus. Tool-parameter defaults and visible key-value history remain candidate carriers that this study does not cover, and the sensitivity result is reported as such --- it was not preregistered and rests on one model. Fourth, the task is a single family: 40 files, one invalidation, one mapping. We do not know how the effect scales with the number of facts that must survive, and the cost multiplier in particular should be expected to depend on how long the agent persists before giving up. Fifth, the factorial decomposition that motivates the redesign survives only as an aggregate: its per-episode carriers were lost when the cloud instance that hosted them was destroyed ahead of its scheduled lifetime, so its discordant counts are cited only where they are derivable from the surviving aggregate, and its uncorrected completion-rate contrasts are withdrawn. Sixth, all three second-iteration models were reached through one gateway with one provider configuration per model, so provider-side behaviour (retry semantics, truncation behaviour, token accounting) is confounded with model identity at the level of absolute token counts; the paired design removes it from the contrast but not from the levels. Seventh, the third iteration is a single-model replication (minimax); its 159-pair estimate is well powered for the preregistered contrast but does not itself establish cross-model robustness, so the direction it suggests is inherited from the three-model second iteration rather than independently re-established there, and its own registered directional prediction was not met. Eighth, the carrier-sensitivity replication stopped at 306 of 400 planned episodes when all three model endpoints hit rate limits in the same window, so its intervals are wider than its design specifies; its post-treatment subgroup conditions on a variable the manipulation itself moves (whether the deletion arm writes before invalidation), which is a selection rather than a controlled comparison; we therefore report the subgroup descriptively in Appendix~\ref{app:sa-strat} and do not treat it as an estimate or a bound, because among the 50 excluded episodes 49 carry the stale-write label and 8 fail to deliver, so the subgroup is selected on the outcome.

\section{Conclusion}

Deleting the entry that records a stale plan did not make the agent fail in the three preregistered iterations; in each, the apparent effect belonged to something else. In the carrier-sensitivity replication it did cost delivery, but only where the deleted content had to be re-fetched. With the prompt allowed to vary, the effect was the wording. With the carrier holding the only copy of the resource identifiers, deletion looked catastrophic because it destroyed the vocabulary for acting. With those identifiers fixed permanently in the prompt and only the version left in the carrier, deletion raises stale-plan writes by 32 points (159 effective pairs, $p=1.2\times10^{-14}$) while delivery is statistically indistinguishable (199 of 200 versus 200 of 200) at essentially equal cost (ratio 1.0177). Moving that version mapping into a retrieval corpus keeps the stale-write effect (79 of 146 effective pairs, +47.3 points) but delivery drops to 135 of 147, so what a deletion costs depends on how hard it is to re-derive the deleted fact, not on whether the fact was stale. The operational lesson is not about staleness but about redundancy: before deleting a memory entry, ask what else still carries each fact it holds. If the fact is stated nowhere else, deletion removes the agent's last route back. If it is stated elsewhere and only a version tag lives in the entry, deletion costs a corrected mis-step and nothing more. If it is stated elsewhere but only expensively retrievable, deletion costs episodes. Deletion is not a memory hygiene primitive with a fixed effect; it is an intervention whose consequence is fixed by whether the deleted entry was the only reachable copy of what the task needed, which is a property of the deployment rather than of the entry.

\appendix

\section{Post-treatment subgroup of the carrier-sensitivity replication}
\label{app:sa-strat}

Fifty of the 147 deletion-arm episodes land no accepted write before the invalidation event at all, so the remaining 97 pairs form a post-treatment subgroup in which the deletion arm did write before invalidation. Conditioning on that subgroup gives 30 of 97 versus 5 for retention, a difference of +25.8 percentage points (exact McNemar $p=4.6\times10^{-6}$) at a token ratio of 1.2452. The subgroup is selected on a variable the manipulation moves, so the number describes that stratum and is not an estimate of the deletion effect; the main text reports the unstratified contrast of +47.3 points over 146 effective pairs.

\begin{table}[t]
\centering
\caption{Post-treatment subgroup of the carrier-sensitivity replication: pairs in which the deletion arm wrote before the invalidation event. Reported descriptively; the subgroup is selected on a variable the manipulation moves.}
\label{tab:6}
\begin{tabular}{lrrrr}
\toprule
Subgroup & $n$ & STALE (del vs keep) & $\Delta$ pp & $\times$ tok \\
\midrule
post-treatment & 97 & 30 vs 5 & +25.8 & $\times$1.2452 \\
\bottomrule
\end{tabular}

\end{table}

\section*{Reproducibility statement}

Every number in this paper is recomputed from recorded judgements and traces by \texttt{recompute\_stats.py} in \texttt{paper\_v2/} for the three-model batches, \texttt{make\_v3\_evidence.py} in \texttt{paper\_v2/scripts/} for the third iteration, and \texttt{summarize\_sa.py} in \texttt{l3\_pilot/} for the carrier-sensitivity replication. Each checks its own expected values against a registered report or a known-gold batch before reading a new one, and each self-test injects a tampered sample size, a reversed direction and an out-of-tolerance deviation and requires each to fail. The instrument's six frozen files are verified by hash against the first iteration's registration; a provenance gate exits non-zero on mismatch. An evidence-chain gate requires each batch's judgements and traces to match its reported record count, and exits non-zero for an aggregate that has lost its carriers. Transport losses are removable only by a tool that keys on the transport flag, with a negative control showing that a genuine failure is not removable. Batch-level reproduction commands, including expected readings, are listed in \texttt{REPRODUCE\_CHECKS.md}; all deviations from the registered plan are in \texttt{DEVIATION\_L3.md}.

\bibliographystyle{plain}
\bibliography{references}
\end{sloppypar}
\end{document}
